\section{Struktur Formal Induksi Matematika Lemah}

Induksi matematika ``lemah'' (\textit{weak induction}) atau disebut juga induksi biasa (\textit{ordinary induction}) merupakan bentuk induksi yang paling umum digunakan. Kata ``lemah'' bukan berarti metode ini kurang kuat, melainkan merujuk pada fakta bahwa langkah induksi hanya mengasumsikan kebenaran $P(k)$ untuk satu nilai $k$ (berbeda dengan induksi kuat yang mengasumsikan seluruh nilai sebelumnya) \cite{rosen2019}.

\subsection{Template Bukti Induksi}

Berikut template standar untuk menyusun bukti induksi:

\begin{enumerate}
    \item \textbf{Pernyataan}: Nyatakan dengan jelas apa yang hendak dibuktikan, misalnya ``Untuk semua $n \geq n_0$, $P(n)$.''

    \item \textbf{Langkah Basis}: Verifikasi $P(n_0)$ secara langsung dengan substitusi $n = n_0$.
    \begin{itemize}
        \item Evaluasi ruas kiri dan ruas kanan secara terpisah.
        \item Tunjukkan bahwa keduanya bernilai sama (atau bahwa pernyataan terpenuhi).
    \end{itemize}

    \item \textbf{Hipotesis Induksi}: Asumsikan $P(k)$ benar untuk sembarang bilangan bulat $k \geq n_0$.

    \item \textbf{Langkah Induksi}: Buktikan $P(k+1)$ benar, menggunakan hipotesis induksi $P(k)$.
    \begin{itemize}
        \item Tuliskan ekspresi $P(k+1)$ secara eksplisit.
        \item Hubungkan dengan $P(k)$ melalui manipulasi aljabar.
        \item Substitusi hipotesis induksi dan sederhanakan.
        \item Tunjukkan bahwa hasilnya memenuhi bentuk $P(k+1)$ yang diinginkan.
    \end{itemize}

    \item \textbf{Kesimpulan}: Berdasarkan prinsip induksi matematika, $P(n)$ terbukti untuk seluruh $n \geq n_0$. $\blacksquare$
\end{enumerate}

\subsection{Contoh Pertama: Jumlah Bilangan Ganjil}

\begin{contoh}
\textbf{Teorema}: Untuk setiap $n \geq 1$, jumlah $n$ bilangan ganjil pertama sama dengan $n^2$:
\[ 1 + 3 + 5 + \ldots + (2n - 1) = n^2 \]

\textbf{Bukti (Induksi):}

\textbf{Langkah Basis} ($n = 1$):
\begin{itemize}
    \item Ruas kiri: suku pertama bilangan ganjil $= 1$.
    \item Ruas kanan: $1^2 = 1$.
    \item $1 = 1$. $\checkmark$
\end{itemize}

\textbf{Hipotesis Induksi}: Asumsikan untuk $n = k$:
\[ 1 + 3 + 5 + \ldots + (2k - 1) = k^2 \]

\textbf{Langkah Induksi}: Buktikan untuk $n = k + 1$:
\[ 1 + 3 + 5 + \ldots + (2k - 1) + (2(k+1) - 1) = (k+1)^2 \]

Ruas kiri:
\begin{align*}
&\underbrace{1 + 3 + 5 + \ldots + (2k - 1)}_{= k^2 \text{ (hipotesis induksi)}} + (2k + 1) \\
&= k^2 + 2k + 1 \\
&= (k + 1)^2
\end{align*}

Ruas kiri $= (k+1)^2 =$ ruas kanan. Langkah induksi terbukti. $\blacksquare$
\end{contoh}

\subsection{Pentingnya Langkah Basis}

\begin{catatan}
Langkah basis \textbf{tidak boleh dilewatkan}, meskipun tampak sepele. Tanpa langkah basis, ``rantai domino'' tidak memiliki titik awal. Langkah induksi hanya mengatakan ``\textit{jika} batu ke-$k$ jatuh, \textit{maka} batu ke-$(k+1)$ juga jatuh'' --- tetapi tanpa memastikan bahwa batu pertama benar-benar jatuh, tidak ada jaminan bahwa prosesnya dimulai.

Contoh berbahaya: tanpa langkah basis, kita bisa ``membuktikan'' bahwa semua bilangan bulat positif sama (yang jelas salah). Ini menunjukkan bahwa kedua langkah (basis \textbf{dan} induksi) secara bersama-sama diperlukan agar bukti valid.
\end{catatan}
